% GENERATED FILE. Edit canonical lessons, then run npm run generate:books.
% canonical-source-hash: fnv1a64:bbbe32c5bbf211fc
% canonical-lessons: KA-C523-samsara, KA-C523-grhini, KA-C523-sangati, KA-C523-viraha, KA-C523-dampatya

\chapter{Partners and Households}
\label{ch:ka-partners-and-households}
\hlchaptermodality{523}

\begin{chapteropening}
\textbf{By the end of this chapter:} \emph{I can say family life, a homemaker, a partner, the pain of parting and married life.}

Complete the last lesson of chapter 523: I can say family life, a homemaker, a partner, the pain of parting and married life.
\end{chapteropening}

\section[\texorpdfstring{\kn{ಸಂಸಾರ} (saṃsāra)}{ಸಂಸಾರ (saṃsāra)}]{\kn{ಸಂಸಾರ} (saṃsāra) — family life; a household}

\label{lesson:KA-C523-samsara}



\begin{quote}
Before the new one: say the Kannada for a sculptor, then the Kannada for to carve.
\end{quote}

\subsection*{You'll want to know: \kn{ಸಂಸಾರ}}
\textbf{\kn{ಸಂಸಾರ}} — \emph{saṃsāra} — \textquotedblleft{}family life; a household\textquotedblright{}.

\begin{quote}
\kn{ಅಕ್ಕ} \kn{ತನ್ನ} \kn{ಸಂಸಾರವನ್ನು} \kn{ಚೆನ್ನಾಗಿ} \kn{ನಡೆಸುತ್ತಾಳೆ}.
\end{quote}

\begin{quote}
\emph{akka tanna saṃsāravannu cennāgi naḍesuttāḷe.}
\end{quote}

\begin{quote}
My elder sister runs her household well.
\end{quote}

\kn{ಸಂಸಾರಿ}, a family man or woman.

\begin{cousinweb}
\textbf{\kn{ಸಂಸಾರ}} (\emph{saṃsāra}) is borrowed from Sanskrit \emph{saṃsāra}, from \emph{saṃ-sṛ}, \textquotedblleft{}to go round\textquotedblright{}, the round of worldly life (Monier-Williams); in Kannada it is the life of a family.
\end{cousinweb}

\begin{grammarlens}[title={What you've built}]
One more word for describing what you lived and felt.
\end{grammarlens}

\subsection*{Your turn}
\begin{itemize}
\raggedright
  \item \emph{Say it:} \emph{saṃsāra}
  \item \emph{Say it:} \emph{saṃsāra}, once more
  \item \emph{Say it:} \emph{akka tanna saṃsāravannu cennāgi naḍesuttāḷe.} — My elder sister runs her household well.
  \item \emph{Recall:} say the Kannada for a sculptor, then the Kannada for to carve, then say \emph{saṃsāra} again
\end{itemize}

\begin{tcolorbox}[breakable,colback=teal!4,colframe=teal!35!black,title={Before you move on}]
What does \kn{ಸಂಸಾರ} mean? (\textquotedblleft{}Family life; a household\textquotedblright{}.) Say it once more.
\end{tcolorbox}
\section[\texorpdfstring{\kn{ಗೃಹಿಣಿ} (gṛhiṇi)}{ಗೃಹಿಣಿ (gṛhiṇi)}]{\kn{ಗೃಹಿಣಿ} (gṛhiṇi) — a homemaker; a housewife}

\label{lesson:KA-C523-grhini}



\begin{quote}
Before the new one: say the Kannada for to carve, then the Kannada for family life.
\end{quote}

\subsection*{You'll want to know: \kn{ಗೃಹಿಣಿ}}
\textbf{\kn{ಗೃಹಿಣಿ}} — \emph{gṛhiṇi} — \textquotedblleft{}a homemaker; a housewife\textquotedblright{}.

\begin{quote}
\kn{ನನ್ನ} \kn{ಅಮ್ಮ} \kn{ಗೃಹಿಣಿ}, \kn{ಮನೆಯ} \kn{ಎಲ್ಲಾ} \kn{ಕೆಲಸ} \kn{ನೋಡಿಕೊಳ್ಳುತ್ತಾರೆ}.
\end{quote}

\begin{quote}
\emph{nanna amma gṛhiṇi, maneya ellā kelasa nōḍikoḷḷuttāre.}
\end{quote}

\begin{quote}
My mother is a homemaker; she looks after all the work of the house.
\end{quote}

\kn{ಗೃಹ}, a house + -\kn{ಇಣಿ}, the woman of it.

\begin{cousinweb}
\textbf{\kn{ಗೃಹಿಣಿ}} (\emph{gṛhiṇi}) is borrowed from Sanskrit \emph{gṛhiṇī}, \textquotedblleft{}the mistress of a house, a wife\textquotedblright{} (Monier-Williams).
\end{cousinweb}

\begin{grammarlens}[title={What you've built}]
One more word for describing what you lived and felt.
\end{grammarlens}

\subsection*{Your turn}
\begin{itemize}
\raggedright
  \item \emph{Say it:} \emph{gṛhiṇi}
  \item \emph{Say it:} \emph{gṛhiṇi}, once more
  \item \emph{Say it:} \emph{nanna amma gṛhiṇi, maneya ellā kelasa nōḍikoḷḷuttāre.} — My mother is a homemaker; she looks after all the work of the house.
  \item \emph{Recall:} say the Kannada for to carve, then the Kannada for family life, then say \emph{gṛhiṇi} again
\end{itemize}

\begin{tcolorbox}[breakable,colback=teal!4,colframe=teal!35!black,title={Before you move on}]
What does \kn{ಗೃಹಿಣಿ} mean? (\textquotedblleft{}A homemaker; a housewife\textquotedblright{}.) Say it once more.
\end{tcolorbox}
\section[\texorpdfstring{\kn{ಸಂಗಾತಿ} (saṅgāti)}{ಸಂಗಾತಿ (saṅgāti)}]{\kn{ಸಂಗಾತಿ} (saṅgāti) — a partner; a companion}

\label{lesson:KA-C523-sangati}



\begin{quote}
Before the new one: say the Kannada for family life, then the Kannada for a homemaker.
\end{quote}

\subsection*{You'll want to know: \kn{ಸಂಗಾತಿ}}
\textbf{\kn{ಸಂಗಾತಿ}} — \emph{saṅgāti} — \textquotedblleft{}a partner; a companion\textquotedblright{}.

\begin{quote}
\kn{ಅಜ್ಜಿ} \kn{ಅಜ್ಜನ} \kn{ಐವತ್ತು} \kn{ವರ್ಷದ} \kn{ಸಂಗಾತಿ}.
\end{quote}

\begin{quote}
\emph{ajji ajjana aivattu varṣada saṅgāti.}
\end{quote}

\begin{quote}
Grandmother has been grandfather's partner for fifty years.
\end{quote}

A life partner, and also a comrade: \kn{ಪ್ರಯಾಣದ} \kn{ಸಂಗಾತಿ}, a travel companion.

\begin{cousinweb}
\textbf{\kn{ಸಂಗಾತಿ}} (\emph{saṅgāti}) is built on Sanskrit \emph{saṅga}, \textquotedblleft{}contact, company\textquotedblright{} (Monier-Williams); the origin of the ending is not settled.
\end{cousinweb}

\begin{grammarlens}[title={What you've built}]
One more word for describing what you lived and felt.
\end{grammarlens}

\subsection*{Your turn}
\begin{itemize}
\raggedright
  \item \emph{Say it:} \emph{saṅgāti}
  \item \emph{Say it:} \emph{saṅgāti}, once more
  \item \emph{Say it:} \emph{ajji ajjana aivattu varṣada saṅgāti.} — Grandmother has been grandfather's partner for fifty years.
  \item \emph{Recall:} say the Kannada for family life, then the Kannada for a homemaker, then say \emph{saṅgāti} again
\end{itemize}

\begin{tcolorbox}[breakable,colback=teal!4,colframe=teal!35!black,title={Before you move on}]
What does \kn{ಸಂಗಾತಿ} mean? (\textquotedblleft{}A partner; a companion\textquotedblright{}.) Say it once more.
\end{tcolorbox}
\section[\texorpdfstring{\kn{ವಿರಹ} (viraha)}{ವಿರಹ (viraha)}]{\kn{ವಿರಹ} (viraha) — the pain of parting; separation from a loved one}

\label{lesson:KA-C523-viraha}



\begin{quote}
Before the new one: say the Kannada for a homemaker, then the Kannada for a partner.
\end{quote}

\subsection*{You'll want to know: \kn{ವಿರಹ}}
\textbf{\kn{ವಿರಹ}} — \emph{viraha} — \textquotedblleft{}the pain of parting; separation from a loved one\textquotedblright{}.

\begin{quote}
\kn{ಗಂಡ} \kn{ವಿದೇಶಕ್ಕೆ} \kn{ಹೋದಾಗ} \kn{ಅವಳಿಗೆ} \kn{ವಿರಹದ} \kn{ನೋವು} \kn{ಆಯಿತು}.
\end{quote}

\begin{quote}
\emph{gaṇḍa vidēśakke hōdāga avaḷige virahada nōvu āyitu.}
\end{quote}

\begin{quote}
When her husband went abroad, she felt the pain of separation.
\end{quote}

A favourite theme of songs and poems.

\begin{cousinweb}
\textbf{\kn{ವಿರಹ}} (\emph{viraha}) is borrowed from Sanskrit \emph{viraha}, \textquotedblleft{}desertion, separation, especially of lovers\textquotedblright{} (CDIAL 11851).
\end{cousinweb}

\begin{grammarlens}[title={What you've built}]
One more word for describing what you lived and felt.
\end{grammarlens}

\subsection*{Your turn}
\begin{itemize}
\raggedright
  \item \emph{Say it:} \emph{viraha}
  \item \emph{Say it:} \emph{viraha}, once more
  \item \emph{Say it:} \emph{gaṇḍa vidēśakke hōdāga avaḷige virahada nōvu āyitu.} — When her husband went abroad, she felt the pain of separation.
  \item \emph{Recall:} say the Kannada for a homemaker, then the Kannada for a partner, then say \emph{viraha} again
\end{itemize}

\begin{tcolorbox}[breakable,colback=teal!4,colframe=teal!35!black,title={Before you move on}]
What does \kn{ವಿರಹ} mean? (\textquotedblleft{}The pain of parting; separation from a loved one\textquotedblright{}.) Say it once more.
\end{tcolorbox}
\section[\texorpdfstring{\kn{ದಾಂಪತ್ಯ} (dāmpatya)}{ದಾಂಪತ್ಯ (dāmpatya)}]{\kn{ದಾಂಪತ್ಯ} (dāmpatya) — married life}

\label{lesson:KA-C523-dampatya}



\begin{quote}
Before the new one: say the Kannada for a partner, then the Kannada for the pain of parting.
\end{quote}

\subsection*{You'll want to know: \kn{ದಾಂಪತ್ಯ}}
\textbf{\kn{ದಾಂಪತ್ಯ}} — \emph{dāmpatya} — \textquotedblleft{}married life\textquotedblright{}.

\begin{quote}
\kn{ಅಜ್ಜ} \kn{ಅಜ್ಜಿಯ} \kn{ದಾಂಪತ್ಯಕ್ಕೆ} \kn{ಐವತ್ತು} \kn{ವರ್ಷ} \kn{ಆಯಿತು}.
\end{quote}

\begin{quote}
\emph{ajja ajjiya dāmpatyakke aivattu varṣa āyitu.}
\end{quote}

\begin{quote}
My grandparents have been married for fifty years.
\end{quote}

From \kn{ದಂಪತಿ}: the life of the pair.

\begin{cousinweb}
\textbf{\kn{ದಾಂಪತ್ಯ}} (\emph{dāmpatya}) is borrowed from Sanskrit \emph{dāmpatya}, \textquotedblleft{}the state of husband and wife, married life\textquotedblright{}, built from \emph{dampati} (Monier-Williams).
\end{cousinweb}

\begin{grammarlens}[title={What you've built}]
One more word for describing what you lived and felt.
\end{grammarlens}

\subsection*{Your turn}
\begin{itemize}
\raggedright
  \item \emph{Say it:} \emph{dāmpatya}
  \item \emph{Say it:} \emph{dāmpatya}, once more
  \item \emph{Say it:} \emph{ajja ajjiya dāmpatyakke aivattu varṣa āyitu.} — My grandparents have been married for fifty years.
  \item \emph{Recall:} say the Kannada for a partner, then the Kannada for the pain of parting, then say \emph{dāmpatya} again
\end{itemize}

\begin{tcolorbox}[breakable,colback=teal!4,colframe=teal!35!black,title={Before you move on}]
What does \kn{ದಾಂಪತ್ಯ} mean? (\textquotedblleft{}Married life\textquotedblright{}.) Say it once more.
\end{tcolorbox}
